% 2007 年浙江卷理科第 8 题：y=f(x) 与 y=f'(x) 在同一坐标系中的四个图象选项。
%   （A）f=x^2 与 f'=2x；（B）f 单调递减而 f' 取正值（不可能）；
%   （C）f=e^x 与 f'=e^x；（D）f 为四次型曲线、f' 在其极值点变号。
% 来源 PDF：gaokaomath-src/普通高考/2007/2007浙江理.pdf
% SHA-256：3b1e5fbca7888b1cff1fc809211aac94f14ea503de0229c4edf7cfb54e5bd2d4
% 原页：2007浙江理.pdf 第 1 页（A3 三栏）第 2 栏，书页码 417。
% 重绘依据：按原页四个小图的曲线形状、相对位置与坐标轴重绘。
\begin{tikzpicture}[scale=0.9, semithick,
  every node/.style={inner sep=1pt, font=\small}]

  % ---------------- （A） ----------------
  \begin{scope}[shift={(0,0)}]
    \draw[->,>=stealth] (-1.8,0) -- (2.1,0) node[right] {$x$};
    \draw[->,>=stealth] (0,-1.4) -- (0,1.8) node[above] {$y$};
    \draw[domain=-1.25:1.25,samples=100] plot (\x,{\x*\x});
    \draw (-1.5,-1.2) -- (1.5,1.2);
    \node[below] at (0.8,-1.7) {（A）};
  \end{scope}

  % ---------------- （B） ----------------
  \begin{scope}[shift={(4.6,0)}]
    \draw[->,>=stealth] (-1.8,0) -- (2.1,0) node[right] {$x$};
    \draw[->,>=stealth] (0,-1.4) -- (0,1.8) node[above] {$y$};
    \draw[domain=-1.4:1.4,samples=100] plot (\x,{1.3*exp(-0.7*(\x+0.2))});
    \draw[domain=-1.4:1.4,samples=100] plot (\x,{0.9*(\x-0.1)});
    \node[below] at (0.8,-1.7) {（B）};
  \end{scope}

  % ---------------- （C） ----------------
  \begin{scope}[shift={(0,-4.4)}]
    \draw[->,>=stealth] (-1.8,0) -- (2.1,0) node[right] {$x$};
    \draw[->,>=stealth] (0,-1.4) -- (0,1.8) node[above] {$y$};
    \draw[domain=-1.3:0.55,samples=100] plot (\x,{0.55*exp(1.1*\x)});
    \draw[domain=-0.2:1.2,samples=100] plot (\x,{0.65*exp(0.8*\x)});
    \node[below] at (0.8,-1.7) {（C）};
  \end{scope}

  % ---------------- （D） ----------------
  \begin{scope}[shift={(4.6,-4.4)}]
    \draw[->,>=stealth] (-1.8,0) -- (2.1,0) node[right] {$x$};
    \draw[->,>=stealth] (0,-1.4) -- (0,1.8) node[above] {$y$};
    \draw[domain=-1.4:1.5,samples=120]
      plot (\x,{0.9*cos(deg(2.2*\x))});
    \draw[domain=-1.4:1.5,samples=120]
      plot (\x,{0.55*sin(deg(2.2*\x))});
    \node[below] at (0.8,-1.7) {（D）};
  \end{scope}
\end{tikzpicture}
